\section{A reserve estimate is not a future payment}
\label{sec:reserving-bootstrap}

The reserve of 148 leaves an important question unanswered: how much might
the insurer ultimately pay? Repeating an estimation procedure can show how
its fitted answer varies. Even if that answer were known exactly, however,
the future claims would still vary. England's retained 2001 working
addendum makes this distinction operational: first resample and refit to
obtain future conditional means, then simulate payments conditional on
those means. The inspected text is Actuarial Research Paper 138; its exact
equivalence to the separately published 2002 journal addendum has not been
established \citep[Sections 2--3 and Appendix A]{england2001addendum}.

Consider a small probability example, separate from the triangle. Suppose
a fitted future mean, denoted by $M$, is either 2 or 8 with equal
probability. These are authored possibilities, not estimates from an
insurer. If $M=2$, let a payment $Y$ be 0 or 4, equally likely;
if $M=8$, let it be 4 or 12, equally likely. Within each fitted case,
the payment has mean $M$ and variance $2M$. Across cases, the fitted mean
has expectation 5 and variance 9. The total variation in payments includes
both the variation left within a case and the variation between case means:
\[
 \operatorname{Var}(Y)
 =\E\!\left[\operatorname{Var}(Y\mid M)\right]
  +\operatorname{Var}\!\left(\E[Y\mid M]\right)
 =\E[2M]+\operatorname{Var}(M)=10+9=19.
\]
One can check the result directly: $Y$ is 0, 4 or 12 with probabilities
$1/4$, $1/2$ and $1/4$. Its mean is 5, and its variance is
$(25/4)+(1/2)+(49/4)=19$. Reporting the distribution of fitted means
alone would instead give variance 9 and no payment above 8. The full
authored payment law gives probability $1/4$ to that event.

Even those two moments do not select a unique payment law. Keep the same
two fitted means, but, conditional on $M=2$, pay 1 or 6 with probabilities
$4/5$ and $1/5$; conditional on $M=8$, pay 6 or 16 with those probabilities.
The conditional means remain 2 and 8. Their variances are
$(4/5)1+(1/5)16=4$ and $(4/5)4+(1/5)64=16$, respectively.
The overall mean and variance remain 5 and 19, but the chance of paying
more than 8 is now $1/10$. Matching a reserve and a prediction error
therefore cannot determine a tail probability or the insurer's compensation
for non-financial risk.

\subsection{What the two stages require}

The addendum starts from a particular model for incremental payments.
For an observed incremental payment $X_{ij}$ in origin period $i$
and development period $j$, its fitted mean $m_{ij}$ is strictly positive.
The model assumes $\E[X_{ij}]=m_{ij}$,
$\operatorname{Var}(X_{ij})=\phi m_{ij}$ and
$\log m_{ij}=c+\alpha_i+\beta_j$, with reference levels fixed to identify
the parameters. This incremental model and Mack's conditional cumulative
model have separate assumptions, even where their fitted reserve amounts
agree. A common point estimate does not make their uncertainty models
interchangeable \citep[equations 2.1--2.3]{england2001addendum}.

The first stage measures each observed increment's departure from its fitted
mean using the Pearson residual $r_{ij}=(X_{ij}-m_{ij})/\sqrt{m_{ij}}$.
With $N$ observed cells and $p$ fitted parameters, the paper estimates
$\widehat\phi=\sum r_{ij}^2/(N-p)$ and resamples residuals after multiplying
them by $\sqrt{N/(N-p)}$. The adjustment must be applied once. It cannot
be followed by another correction of the same estimation variance
\citep[equations 2.4--2.6 and 3.1]{england2001addendum}.

Start from the observed latest cumulative amounts, work backwards through
fitted development factors, and difference the fitted cumulative amounts
into increments. Calculate the residuals and dispersion from those increments.
For each bootstrap draw, resample adjusted residuals with replacement and
invert the residual relation to form pseudo-increments. Accumulate them,
refit the development factors, project future cumulative claims and difference
them into future means. Retain the same fitted draw across all cells. This
produces a joint array of conditional means, not yet a joint array of payments
\citep[Section 2 and Appendix A, Triangles 1--9]{england2001addendum}.

The second stage chooses a process law. For a positive future mean $m$
and positive dispersion $\phi$, the paper proposes $\phi Z$, where $Z$
has a Poisson law with mean $m/\phi$. Its mean is $m$ and variance is
$\phi m$. A gamma law with shape $m/\phi$ and scale $\phi$ has the
same two moments but different support and tail. That is a modeling choice,
not a change of random-number syntax. Each future payment is drawn
conditionally, and aligned payments are summed by origin, payment period
and total \citep[Section 3 and Appendix A, Triangles 9--10]{england2001addendum}.

Dependence must survive aggregation. In the first small example, suppose
two cells share the same $M$ but are independent conditional on it. Their
conditional total variance is $4M$, with expectation 20, while their
conditional total mean is $2M$, with variance 36. Their total variance
is therefore 56. Independently redrawing each cell from its marginal
payment law would give $19+19=38$, losing the dependence created by the
shared fit. Marginal distributions are insufficient instructions for a
joint simulation.

\subsection{Sparse data, failed draws and implementation choices}

The six observed increments in the earlier triangle fit the incremental
model exactly. It has $p=2(3)-1=5$ parameters and only $N-p=1$ residual
degree of freedom, giving zero fitted dispersion. This arithmetic does
not establish that future payments are deterministic. It also does not
identify Mack's final development variance. The proposed interface returns
\code{DISPERSION\_UNRESOLVED} for this illustrative input and retains the
point reserve of 148 separately. More data or a justified additional
assumption is required before a predictive distribution can be admitted.

Every draw records its model, source version, data vintage, residual pool,
fitted parameters, process law and random stream. Nonpositive degrees of
freedom, an unsupported negative mean, a zero projection denominator or
non-finite arithmetic receives an explicit failure. Structurally absent
triangle cells have a separate mask. A failed payment is not converted
to zero or silently replaced by another draw. In the authored example,
if the $M=8$ case cannot be evaluated, its probability $1/2$ remains
unresolved: the known contribution to the full mean is $\tfrac12(2)=1$,
and a complete mean is withheld. In a finite bootstrap, attempted
failure frequency is reported as such; it does not identify population
probability.

The retained source itself needs a qualification. Its negative-mean
suggestion draws $X$ with mean $|m|$ and then subtracts $2m$. Read
literally for $m<0$, the resulting mean is $|m|-2m=-3m$, not $m$;
at $m=-4$ it is 12. That prescription is wrong relative to its stated
mean target in the inspected working version. A translation or reflection
can preserve a negative mean, but changes the support and requires a
different variance model. AccDSL leaves that extension unsupported
pending an explicit signed-claims model; this local finding does not
reject the positive-mean two-stage procedure
\citep[Section 3, final paragraph]{england2001addendum}.

Software names also conceal consequential choices. In the inspected
ChainLadder 0.2.17 source snapshot, \code{od.pois} calls a negative-binomial
generator for dispersion greater than one. The paper's scaled Poisson
and that negative binomial can share a mean and variance while assigning
different probabilities to possible payments. The same snapshot uses
signed gamma draws and replaces missing process values by zero. An adapter
must distinguish structural absence from failed computation, identify the
exact law and its parameter domain, and preserve separate resampling and
process streams before its output can satisfy this proposal. These are
source-code observations, not completed runtime validation
\citep{chainladder2023snapshot}.

The paper compares 1,000 simulations with analytic and earlier bootstrap
summaries. That provides a worked implementation comparison under its
assumptions. It does not establish future tail coverage, select an insurer's
risk adjustment or justify 1,000 draws as a universal budget. A later
empirical protocol must address calendar effects, tail development, pooling,
heldout diagonals, simulation uncertainty and consequences for measured
insurance groups \citep[Section 4, Tables 1--3 and Figure 1]{england2001addendum}.

\subsection{The measurement date changes the question}

A lifetime payment distribution differs from a one-year change in
estimated liabilities. The 2019 accepted manuscript by England, Verrall
and W\"uthrich defines the claims development result as opening reserves
less payments during the year and less the revised closing reserve. The
closing estimate uses the information that has emerged during that year
\citep[equations 5--7 and Section 3.2]{england2019lifetime}.
If the opening reserve is 100, payments are 30 and the revised closing
reserve is 80, the result is $100-30-80=-10$. Subtracting cash alone
would yield 70 and answer a different question. The scenario must carry
both realized next-year payments and the prescribed re-estimation rule,
with the correct information cutoff. A lifetime draw cannot be relabeled
as a one-year result. Neither statistic by itself supplies the discounting,
group allocation or risk compensation required for insurance measurement.
